\subsection{Experimental Design}\label{sec:configurations}
This section defines what is trained, including the cost function branches, the network and training configurations, and the number of seeds per run. The full grid of results is given in App.~\ref{app:grid}. As described in Sec.~\ref{sec:fin-gan}, the eight cost function combinations of \textcite{vuleticFinGANForecastingClassifying2024} were used:
$\mathrm{PnL}^{*}_{a}$;
$\mathrm{SR}^{*}$;
$\mathrm{PnL}^{*}_{a}$ \& $\mathrm{MSE}$;
$\mathrm{PnL}^{*}_{a}$ \& $\mathrm{SR}^{*}$;
$\mathrm{PnL}^{*}_{a}$ \& $\mathrm{STD}$;
$\mathrm{SR}^{*}$ \& $\mathrm{MSE}$;
$\mathrm{PnL}^{*}_{a}$ \& $\mathrm{MSE}$ \& $\mathrm{SR}^{*}$;
$\mathrm{PnL}^{*}_{a}$ \& $\mathrm{MSE}$ \& $\mathrm{STD}$. 

A ninth combination, using only the $\mathrm{MSE}$, was added as a supervised counterpart to ForGAN. A tenth branch was trained on cross-entropy alone, giving the ForGAN baseline. The ten branches are referred to by their cost function terms throughout Chapter~\ref{sec:exp}, with the nine Fin-GAN branches named after their economics-driven terms. Each combination was trained under several configurations, listed in Table~\ref{tab:configurations}. 

\begin{table}[H]
\centering
\begin{tabular}{lccccc}
\toprule
Configuration & $l$ & $N$ & $R_G$ & $R_D$ & $n_{\text{grad}} + n_{\text{epochs}}$ \\
\midrule
1  & 10 & 8  & 8  & 8  & $25 + 100$   \\
2  & 10 & 8  & 8  & 8  & $100 + 500$  \\
3  & 10 & 32 & 8  & 8  & $100 + 500$  \\
4  & 10 & 8  & 8  & 64 & $100 + 500$  \\
5  & 10 & 32 & 8  & 64 & $100 + 500$  \\
6  & 30 & 32 & 8  & 64 & $100 + 500$  \\
7  & 10 & 32 & 8  & 64 & $100 + 1500$ \\
8  & 10 & 32 & 64 & 64 & $100 + 500$  \\
9  & 10 & 32 & 64 & 64 & $100 + 1500$ \\
10 & 10 & 8  & 64 & 64 & $100 + 500$  \\
\bottomrule
\end{tabular}
\caption{Model configurations}
\label{tab:configurations}
\end{table}

The configurations differ in one hyperparameter at a time, to isolate the effect of each change. The first configuration is the reference configuration, following the values and the schedule of Algorithm 1 in \textcite{vuleticFinGANForecastingClassifying2024}. The second configuration tests whether a longer training period improves the distributional quality. The third and fourth configurations increase $N$ and $R_D$ separately. 

The fifth configuration increases both, matching the default values of the published implementation. The sixth configuration tests whether a longer condition window of \num{30} trading days carries more information about the volatility of the series. The seventh configuration extends the training further. The eighth configuration increases the generator hidden size $R_G$, which is not varied in the paper or in the published implementation. The ninth configuration combines the longer training with the increased $R_G$. Lastly, the tenth configuration combines the increased $R_G$ with the original $N$. An additional test run was performed with the scaling parameters computed from the entire training split, using configuration \num{5} of Table~\ref{tab:configurations}.

Each configuration was trained with five fixed random seeds. Training runs of the same configuration differ noticeably in their distributional statistics, due to the stochastic nature of the training process. Therefore, all results were reported as mean and seed spread. Fixing the seeds makes the runs reproducible in the environment described in Sec.~\ref{sec:software}.
